\section{How Common Are Deferred Egress Paths?}
\label{sec:measure}

Before building a mechanism, we asked how often apps combine the three ingredients of the problem: persisted deferred
work, a network sink reachable from a deferred-job entry point, and a user-facing telemetry or consent control. This is a
\emph{candidate} screen. Static analysis cannot see runtime state or app wording, so it cannot establish a violation.

\paragraph{Corpus.} From the F-Droid index~\cite{fdroid} (fetched 2026-11-08, SHA-256 recorded), we took the latest version
of every app that requests \texttt{INTERNET} and is at most 80\,MB, giving \eligibleN{} apps. We selected all \trackN{} that
carry F-Droid's \emph{Tracking} anti-feature, plus a seeded random sample of \randN{} others. All \corpusN{} APKs were
downloaded from \texttt{f-droid.org} and verified against the index hashes. F-Droid apps are open source and often minified,
so library method names may be renamed. We therefore detect WorkManager through \texttt{ListenableWorker} subclasses,
whose names are kept for reflection, as well as through API references.

\paragraph{Analysis.} Our scanner (\S\ref{sec:impl}) finds deferred-job entry points: \texttt{doWork}/\texttt{startWork} of
WorkManager workers, \texttt{onStartJob} of \texttt{JobService}s, and \texttt{onPerformSync}. From each, it searches
breadth-first for a network sink (\texttt{HttpURLConnection}, OkHttp, Retrofit, Volley, Ktor, raw sockets and similar). The
search is over-approximate. We validated it on a seeded sample of \valN{} app-code entries in apps with a consent signal,
judging from class names and app purpose whether network use is on the job's own path. Of \valShallowN{} entries within three hops, \valShallowYes{} were judged plausible and \valShallowUnclear{} unclear;
of \valDeepN{} deeper entries, \valDeepYes{} were plausible, \valDeepUnclear{} unclear and \valDeepNo{} implausible. These class-name/app-purpose
judgments do not validate runtime paths or estimate false negatives. We report two screening thresholds
(three hops and any depth), not lower and upper bounds on real flows. Consent signals are identifier-like constants pairing a topic token
(\textit{analytics}, \textit{telemetry}, \textit{crash}, \textit{consent}, ...) with a control token (\textit{enable},
\textit{opt}, \textit{send}, ...), or English toggle-like resource strings such as ``Share anonymous usage statistics''.

\begin{table}[t]
\centering
\caption{F-Droid corpus screen: apps (\%). Random-stratum percentages carry Wilson 95\% intervals conditional on the screen. Tracking is a census of eligible tracking apps. $^\dagger$persisted deferred work + entry$\to$network + consent signal.}
\label{tab:corpus}
\mtab
\begin{tabular}{@{}lrr@{}}
\toprule
\mhead Feature & Tracking ($n{=}$60) & Random ($n{=}$300) \\
\midrule
Network API & 57 (95) & 277 (92) [89, 95] \\
\texttt{JobScheduler} & 32 (53) & 183 (61) [55, 66] \\
\quad \texttt{setPersisted} & 22 (37) & 114 (38) [33, 44] \\
WorkManager & 9 (15) & 45 (15) [11, 19] \\
AlarmManager & 37 (62) & 188 (63) [57, 68] \\
SyncAdapter & 0 (0) & 3 (1) [0, 3] \\
Persisted deferred work (any) & 22 (37) & 116 (39) [33, 44] \\
Entry $\to$ network, $\le$3 hops & 5 (8) & 17 (6) [4, 9] \\
Entry $\to$ network, any depth & 13 (22) & 53 (18) [14, 22] \\
Persisted + entry$\to$net, $\le$3 & 4 (7) & 17 (6) [4, 9] \\
Persisted + entry$\to$net, any & 10 (17) & 50 (17) [13, 21] \\
Consent/telemetry signal & 23 (38) & 42 (14) [11, 18] \\
Candidate, $\le$3 hops$^\dagger$ & 3 (5) & 5 (2) [1, 4] \\
Candidate, any depth$^\dagger$ & 6 (10) & 12 (4) [2, 7] \\
Telemetry/crash SDK & 19 (32) & 14 (5) [3, 8] \\
Native libraries & 38 (63) & 210 (70) [65, 75] \\
\bottomrule
\end{tabular}
\end{table}


\paragraph{Findings.} Table~\ref{tab:corpus} reports each sampling stratum separately: the tracking stratum is a census
within the eligible index, random-stratum intervals are conditional on classification, and we do not pool the
oversampled strata. With only \valN{} validation judgments, the plausible fraction is \valYes{}/\valN{} (Wilson 95\%
interval from \valYesLo{}\% to \valYesHi{}\%), and \valUnclear{} remain unclear. Summed over both strata (a count of
apps, not a pooled rate), between \candConsentLoN{} (three-hop screen) and \candConsentHiN{} (any depth) screened apps persist deferred work, reach the network from a deferred-job entry
point and carry a consent or telemetry signal; neither count estimates a contract-violation prevalence. Only
\pctPollStop{}\% of app-code workers poll \texttt{isStopped()}, which does not show whether cancellation interrupts
uploads; between-batch stopping is a fixture assumption, not a corpus finding. Native code, which most screened apps in
both strata ship, is a path no Java-level fence mediates.

\paragraph{Stated promises.} A consent signal does not say what an app promises. We coded the user-facing wording of
all \promiseApps{} apps with a consent signal (\promiseCand{} of them candidates), using the toggle-like strings in each APK
and its F-Droid listing; privacy policies were not collected. A language model applied written rules to every app and
flagged \promiseFlagged{} for review; an author re-read those and recoded \promiseRecoded{}. In this sample,
\promisePending{} apps word a switch as sending or sharing data (\promiseCandPending{} of the candidates), which we read as
also covering queued data; \promiseCollectOnly{} word it as collection only; \promisePrompt{} only offer to send an
individual report; \promiseUnclear{} expose only code identifiers or unclear wording; \promiseFalse{} are scanner false
positives. \promiseDelete{} offer deletion and \promiseRestrict{} a purpose restriction. Only the first group makes
queued egress after withdrawal a contract violation. These counts describe the surveyed apps, not F-Droid.

ACRA appears in \acraApps{} apps, but its preference keys are library constants and do not show an exposed toggle
(\S\ref{sec:eval-acra}). The screen shows that the ingredients co-occur in the screened apps, not how often they combine into a violation;
the rest of the paper turns to controlled, contract-aware experiments.
